% Owner-confirmed tracking and re-identification do not imply sports-production or AWS experience.
\renewcommand{\cvfootertitle}{Cover Letter}
\renewcommand{\cvBody}{%
    \cvletterhead
        {Stuttgart, 25 October 2026}
        {SkillCorner\\Computer Vision team\\Paris, France}
        {Application: Computer Vision Scientist}

    Dear SkillCorner team,

    % Earlier opening retained: Turning broadcast video into reliable player-level data requires more than a strong model: it requires a pipeline whose predictions remain useful outside the training set.
    While studying computer science, I played American football and found myself imagining exactly the kinds of sports-analysis applications SkillCorner develops today. As a player, I would have found those tools immensely useful. The opportunity to help build them now, combining my experience on the field with my work in computer vision, makes this role particularly compelling.

    % Earlier project evidence retained: RemaNet activity-recognition camera pipelines used ONNX inference and PyTorch training on NVIDIA GPUs; the industrial safety application used four cameras with batched inference.
    Your requirement to design, train and optimize deep learning models for detection, tracking, re-identification and recognition closely matches what I did during my first six months at Fraunhofer IPA. Since joining, I have owned the full ML lifecycle of more than four industrial vision projects, from dataset creation and annotation through PyTorch training, evaluation and deployment. My work includes RemaNet's activity-recognition camera pipelines and person detection across four cameras with batched inference. I also deployed INT8-quantized object detection at approximately 14 frames per second on a Berghof PLC fitted with a Hailo accelerator, iterating with customers on field feedback.

    % Earlier transition retained: I am especially interested in extending my detection and activity-recognition experience into multi-object tracking and re-identification on broadcast video.
    Evaluation is a particular research interest of mine. I co-authored \textit{Semantic Robustness Probing via Inpainting}, an ECML PKDD 2026 demo that uses diffusion-based inpainting to probe object detectors against operating-condition factors. At SkillCorner, I would bring that investigative approach to understanding where predictions fail and which improvements matter downstream, applying my experience to the challenges of broadcast sports video.

    I bring the engineering habits needed to turn experiments into maintainable systems. I maintain institute repositories in Python, Rust and C++, work with GitLab CI/CD, and use MLflow and ClearML for experiment tracking. I review agent-generated changes and run tests, linters and builds before accepting them. My advanced university coursework in AI and computer graphics, together with my master's degree in computer science, also gives me a foundation for discussing the theory and geometry behind implementation choices.

    % Earlier personal connection retained: Outside work, I enjoy cycling; professionally, I am motivated by the opportunity to make vision research useful to people making real decisions.
    Having played the sport myself, I appreciate why reliable tracking data can matter to players and coaches. I would welcome the chance to help turn research improvements into tools they can use. I work in English at C1 level, speak French at B1 and am open to relocating to Paris.

    \cvletterclosing{Best regards}
}
